\chapter[An accountable ensemble with bounded resolution]{An accountable ensemble\\with bounded resolution}
\label{ch:ensemble}

% BEGIN SOURCE UNIT M06:00
\section{The design objective and its limits}
\label{merge:M06:00}


The proposed system must make a dangerous path difficult: a plausible reading is
accepted, translated correctly, and released without anyone discovering that it
omits a material qualification. Multiple generators alone do not close that path.
They may share a source extraction, a vocabulary, a training distribution and the
same mistaken assumption. The design therefore uses redundancy across different
activities, not just repeated answers to one question.

One activity inventories the source and its dependencies. Another identifies
normative statements. Several candidate constructors express different readings.
An opponent seeks a missing condition or counterexample. A formal evaluator
checks consequences. A separate reference process supplies expected decisions
for reviewed scenarios. A Java verifier checks preservation. Each activity has a
defined responsibility and a defined way to fail. A failure or discrepancy is
recorded even when every other member appears successful.

The main output is an interpretation dossier. It contains the proposed control,
the source it covers, the readings considered, the reasons for their dispositions,
the tests and formal checks, and the uncertainty that remains. The controller
may prepare this dossier for review. It does not acquire authority to approve
legal meaning merely because it has completed its work.

There are three success conditions at different levels. An engineering run
succeeds when it follows the specified state machine and produces a complete
report within its limits. An interpretation investigation succeeds when it
resolves its declared questions with adequate evidence or precisely identifies
what requires judgment. A release succeeds only after the required people approve
the exact interpretation and executable package. A run can satisfy the first
condition while ending with an unresolved legal question. That is a valid and
useful outcome, not an excuse to fabricate certainty.

No finite ensemble can establish that it has found every possible reading of
unrestricted language. The operational objective is more precise: cover a
reviewed source perimeter; require explicit treatment of every inventoried unit;
search a declared set of interpretive dimensions; investigate every detected
discrepancy under finite limits; and prevent a material unresolved question from
being silently converted into release. These obligations can be implemented and
tested. Their empirical effectiveness at discovering unknown omissions requires
the evaluation in Chapter~\ref{ch:evaluation}.
% END SOURCE UNIT M06:00

% BEGIN SOURCE UNIT M06:01
\section{Why diversity has to be designed}
\label{merge:M06:01}


Repeated sampling can reduce some random mistakes while leaving a shared mistake
untouched. A simple probability model makes the distinction clear. Let $q$ be the
probability of a common failure that makes every member wrong. Conditional on
that failure not occurring, assume for this illustration that each of $m$ members
has independent error probability $p$. The probability that all members are wrong
is then
\begin{equation}
 \Pr(\text{all wrong})=q+(1-q)p^m.
 \label{eq:common-error}
\end{equation}
The two events are disjoint: either the common failure occurs, or it does not and
all residual errors occur together. This derivation is an illustrative model,
not an estimated error law for language models or lawyers.

If $q=0.02$, $p=0.1$ and $m=3$, the expression gives $0.02098$. Treating the three
members as unconditionally independent would instead give $0.001$. Adding
members reduces the residual term in this example, but the common failure
remains. A missing footnote in the shared extraction is a concrete way that the
common failure could arise. Three careful interpretations of the same incomplete
text may agree perfectly.

A second calculation concerns the variability of a mean of member indicators.
Suppose $X_i$ records whether member $i$ makes a particular error, all have variance
$\sigma^2$, and every pair has correlation $\rho$. Expanding the variance of their
mean gives
\begin{align}
 \operatorname{Var}\!\left(\frac1m\sum_{i=1}^m X_i\right)
 &=\frac{1}{m^2}\left(m\sigma^2+m(m-1)\rho\sigma^2\right)\notag\\
 &=\frac{\sigma^2}{m}\{1+(m-1)\rho\}.
 \label{eq:ensemble-variance}
\end{align}
The expression explains why counting correlated members exaggerates the amount
of independent information. It does not identify their correlation in practice,
and it does not supply a legal error probability. Its role is to motivate
measurement of shared failure and to discourage confidence claims based on panel
size alone.

Diversity should therefore be recorded along several axes. Members may use
different model families, prompts, representations, source acquisition routes,
retrieval methods, factual schemas or verification methods. These differences are
opportunities for complementary errors; they do not prove independence. Two
vendors can share training material. Two parsers can both omit an image. Two human
reviewers can both receive a leading case summary. The record should describe
what differs and what remains shared.

The strongest inexpensive separation is often procedural. Let source inventory
and initial interpretation happen before either sees the other's preferred
conclusion. Freeze the initial outputs. Only then expose disagreements for
investigation. This preserves evidence that one member independently noticed a
condition. If all members see the first answer before responding, later agreement
may merely reflect anchoring. Deliberation can still be useful, but it must not
rewrite the history of independent proposals.
% END SOURCE UNIT M06:01

% BEGIN SOURCE UNIT M06:02
\section{The ensemble members and their responsibilities}
\label{merge:M06:02}


The prototype uses roles rather than assuming that every role is a language
model. The source custodian acquires the official document, retains the exact
bytes and metadata, and records incorporated materials. A document inventory
process enumerates paragraphs, footnotes, tables, attachments and visible
cross-references. A second inventory method compares the enumeration with the
rendered document. Both methods may be automated initially, but unexplained
mismatches remain review questions.

The normative reader identifies who must or may do what, under which conditions,
with which exceptions and dates. A controlled-language reader expresses the same
material using a restricted template. Their outputs are compared at the level
of actor, activity, modality, object, condition, exception, time and dependency.
This comparison can reveal a missing phrase even when the two prose summaries
sound similar.

A counterinterpretation reader is required to propose consequential alternatives
where the source supports them, and to say when it found none. It is not required
to invent a controversy for every sentence. Its task includes seeking an omitted
qualification, an alternative attachment of an exception, or a classification
that changes a boundary case. Unsupported alternatives are preserved as such,
which helps distinguish deliberate testing from legal disagreement.

The formal checker validates types, references and explicit semantics, compares
candidate rules over supported domains, and replays witnesses. It has no
permission to promote a source assumption to a fact. The scenario designer
constructs examples that separate readings and includes unknown, conflicting and
out-of-scope evidence. Its expected outcomes come from an independent review
process or remain labelled provisional.

The implementation checker compares the accepted rule with the Java package and
verifies the package's exact identity. Finally, the dossier assembler reconciles
all required outputs and creates the review report. It may summarize, but it may
not suppress a minority objection, change its severity or mark it resolved without
a linked resolution record.

A role can fail in several ways: unavailable, malformed output, unsupported
claim, resource limit, incompatible source version or abstention. Abstention is
an informative response when a role lacks the evidence needed for its question.
It is not a vote for the majority. If a role is mandatory for the configured
coverage claim, its missing output prevents that claim from being marked complete.
A replacement role must be explicitly authorized by the configuration and its
independence limits recorded.

\begin{table}[htbp]
\centering\small
\caption{Responsibilities that create different opportunities to detect error.}
\label{tab:ensemble-roles}
\begin{tabularx}{\textwidth}{P{31mm}YY}
\toprule
Role & Question answered & Failure that must remain visible\\
\midrule
Source inventory & What material is present and incorporated? & Missing attachment, unmatched paragraph, inaccessible dependency\\
Normative reader & What obligations, conditions and exceptions are expressed? & Omitted modality, scope or qualification\\
Alternative reader & What supported choice changes a consequence? & No supported alternative found; unexplored family\\
Argument critic & What defeats the premises or inference? & Unanswered objection or unverified authority\\
Formal checker & What follows from this specified rule? & Invalid encoding, unknown solver result, limited domain\\
Reference reviewer & What is expected for this scenario, and why? & Disagreement, incomplete facts, shared authorship\\
Java checker & Does this package preserve the accepted specification? & Mismatch, unsupported operation, stale build\\
Dossier assembler & Are all required results and dispositions present? & Missing member, open issue, broken evidence link\\
\bottomrule
\end{tabularx}
\end{table}
% END SOURCE UNIT M06:02

% BEGIN SOURCE UNIT M06:03
\section{Build a source perimeter before interpreting it}
\label{merge:M06:03}


The source perimeter is the set of documents and versions the investigation
claims to cover. It begins with the named circular, but incorporated definitions,
annexes and transition provisions may expand it. Each dependency has a reason:
explicit incorporation, definition of a term, amendment, background explanation,
or a proposed analogy. These relationships must not all be collapsed into a
generic ``relevant document'' link.

For 23EC46, the retained inventory has eighteen paragraphs and four footnotes.
The executable slice covers paragraph 10. That is a legitimate prototype scope
only if the report makes the remaining material visible and explains whether it
is contextual, another obligation, a dependency, or outside the selected slice.
Labelling all unimplemented paragraphs ``not applicable'' merely because the
prototype lacks code would be wrong. Implementation capability is not a legal
applicability test.

The inventory assigns each source unit a stable locator within a particular
source hash. It records the extracted text, its position, and whether the unit
contains a possible normative statement or reference. The second inventory can
use a different parser or rendered-page inspection. A mismatch creates a source
coverage discrepancy before any interpretation ranking takes place.

The controller expands dependencies until each required edge reaches a retained
source, a justified out-of-scope disposition, or an unresolved dependency record.
Cycles are legal as document relationships: a code and a guidance note can refer
to each other. The traversal uses visited version identifiers to avoid infinite
retrieval, while preserving both edges. A retrieval depth limit stops processing
but leaves the unexpanded edge visible. It does not prove that the dependency is
unimportant.

Publication date, effective date and retrieval date occupy separate fields. A
circular published today may describe a later transition or an earlier
arrangement. The source custodian cannot infer commencement solely from the HTTP
response time or the first date in the text. The retained source packet is a
historical evidence object. A current applicability assessment requires a
separate search for amendments and supersession, with its own cut-off date.

This inventory is an important defense against common omissions, but it is not
an oracle. An extractor might enumerate a footnote without understanding its
legal effect. The requirement that each unit receives a reasoned disposition
connects inventory to interpretation. A coverage check should fail if a unit is
present but has no explanation, or if an explanation claims executable coverage
without identifying the rule and verification evidence that provide it.
% END SOURCE UNIT M06:03

% BEGIN SOURCE UNIT M06:04
\section{Freeze blind proposals before reconciliation}
\label{merge:M06:04}


An initial proposal request contains the frozen source packet, the selected
activity profile, the vocabulary already approved for that profile, and the
member's task. It excludes other members' proposed conclusions. The response is
stored with model or tool identity, configuration, prompt-template version, source
hashes and completion status. Raw output is retained subject to the bank's access
and retention policy; the structured candidate is a validated derivative.

A candidate cannot refer to a source version it was not given without opening a
new retrieval request. This rule prevents a model's recalled description of an
unseen amendment from entering as if it were inspected evidence. Recalled
material can nominate a search, but the supporting claim remains unverified
until the source is acquired and checked.

The blind phase ends when every required member has responded or reached a
terminal failure under the initial-phase budget. The controller records that
boundary. It then compares the candidates and inventories. Later repairs may
share all relevant evidence, because their purpose is informed resolution. The
record must continue to distinguish the initial independent proposals from
responses influenced by the shared discussion.

This separation prevents a misleading report such as ``five independent members
agreed'' when four were shown the first member's answer. It also makes a useful
experiment possible: measure which defects were caught before deliberation and
which required shared investigation. A system that produces agreement by
suppressing dissent should look different from one that resolves an objection
with a missing official definition.

Candidate revisions are immutable. A revised rule points to its parent, states
what changed, and links the evidence or diagnostic that motivated the change.
The original remains available. Reviewer approvals attach to a particular
candidate hash, not to a mutable display name such as ``latest rule.'' If a
source, assumption or rule changes, the old approval cannot silently authorize
the new candidate.
% END SOURCE UNIT M06:04

% BEGIN SOURCE UNIT M06:05
\section{Detect discrepancies at the right level}
\label{merge:M06:05}


A discrepancy is a structured difference that may affect coverage, meaning,
execution or evidence. It need not be a pair of contradictory prose answers.
One member may omit an input that another requires; two members may agree on a
rule but cite different versions; one parser may miss a footnote; the Java
runtime may return the same decision with an incomplete audit trace. All deserve
a recorded disposition appropriate to their consequence.

The comparator first checks identity: source versions, activity profile, subject
unit, candidate dependencies and evaluation times. Comparing answers to different
questions creates false disagreement. A model interpreting a whole offer and one
interpreting a single benefit must be identified as differing in subject identity
before their truth values are compared.

It then compares structured meaning: actor, action, object, modality, conditions,
exceptions, quantifiers, temporal boundaries and references. Exact structural
matches are cheap evidence of agreement about those fields, but source coverage
still needs a separate check. Differences in presentation can be classified as
nonmaterial only with an explicit rationale, such as identical canonical syntax
or a domain-qualified equivalence result.

Formal outcome comparison comes next. A counterexample is especially useful when
two apparently similar rules disagree only at a boundary. For paragraph 10,
omitting product-type linkage is exposed by $P=\False$, $T=\True$, with the other
conditions established. Omitting the discount exception is exposed by a genuine
fee discount. Neither example resolves whether a particular cashback is a fee
discount; that is a different discrepancy with a different evidence need.

The comparator also checks omissions against independent inventories. Every
required source unit, dependency and ensemble role must have a disposition. An
all-agreeing set of candidates can therefore still produce discrepancies. This
is the intended behavior: agreement among candidates cannot erase an uncovered
paragraph or a failed source acquisition.

A discrepancy record includes its origin, affected source units, candidate
identifiers, issue family, factual or legal question, observed difference,
materiality, evidence needed, permitted next actions and current status. It also
stores the consequence of leaving it unresolved. A classification issue might
block a particular incentive category; a corrupted source hash might invalidate
the entire run. The scope of the block must be explicit and conservative when
impact cannot yet be determined.
% END SOURCE UNIT M06:05

% BEGIN SOURCE UNIT M06:06
\section{An issue needs an identity that survives retries}
\label{merge:M06:06}


Naive retry systems often attach the attempt counter to a model response. When
the response creates a revised candidate, a new counter starts. A difficult issue
can then consume unlimited attempts by repeatedly changing its wording or
creating descendants. The proposed controller attaches budgets to the run and to
an issue root that survives candidate revisions.

An issue root identifies the source context, the substantive question and the
affected subject or rule family. It does not include a candidate hash as the only
identity, because candidate hashes necessarily change during repair. The system
assigns an immutable root identifier when the question is first recorded.
Subsequent discrepancies can link to that root or create a child issue. A child
inherits the root's budget account and the run's limits. Creating a new display
record never creates new spending authority.

Automatic deduplication proposes that two issues have the same root when they
concern the same source unit, semantic dimension and consequence. A reviewer or
strict deterministic criterion must resolve uncertain merges. An erroneous merge
could hide a second problem; an erroneous split could waste work. Global action
limits protect against unlimited spending even when deduplication is imperfect,
while the report shows possible duplicates for later review.

Issue identity also separates changed evidence from repeated reasoning. If an
unavailable FAQ becomes available during the same run, the controller can attach
the new version and continue within the remaining budget. If a later circular
changes the applicable rule, that is a new source context and normally a new run.
The old run remains closed with its original result. Opening the new run requires
a recorded trigger and authorization; a worker cannot evade exhausted limits by
automatically renaming the investigation.

A candidate that depends on several issues keeps all of their identifiers. A
resolution of one issue cannot clear the others. The dependency relation also
supports invalidation: retracting a definition affects every candidate that used
it, even if those candidates were reached through different search branches.
The service must query this relation directly rather than trying to reconstruct
it from narrative summaries.
% END SOURCE UNIT M06:06

% BEGIN SOURCE UNIT M06:07
\section{Statuses describe evidence, not optimism}
\label{merge:M06:07}


The issue lifecycle begins at \texttt{OPEN}. An issued action moves it to
\texttt{INVESTIGATING}. It may then return to \texttt{OPEN} with additional evidence,
move to \texttt{AWAITING\_EVIDENCE}, or require \texttt{HUMAN\_REQUIRED}. These are
unresolved states. They differ in the next useful action, not in whether the
system may assume the preferred answer.

Resolution has three main forms. \texttt{RESOLVED\_EVIDENCE} means that an inspected
source or factual record establishes the missing premise under the accepted
review rules. \texttt{RESOLVED\_ADJUDICATION} means that an authorized reviewer has
made and explained the interpretive decision. \texttt{RESOLVED\_EQUIVALENCE} means
that the apparently different executions have been shown equivalent over a
specified domain and that no distinct source commitment remains at issue.
These statuses require different evidence and cannot substitute for one another.

For example, an SMT equivalence result cannot resolve a dispute about whether the
word ``may'' establishes a universal classification. Both candidates may have
been encoded with the same unjustified classification flag. Conversely, an
explicit source passage that defeats a missing-exception candidate need not wait
for a probabilistic vote. The source-based resolution identifies the clause,
updates the candidate and requires the affected formal and Java checks again.

An issue may become terminal while remaining unresolved. The terminal reason can
be \texttt{BUDGET\_EXHAUSTED}, \texttt{DEADLINE\_REACHED}, \texttt{NO\_PROGRESS},
\texttt{DEPENDENCY\_UNAVAILABLE}, \texttt{CANCELLED}, or an integrity failure.
Terminality means automatic work has stopped. It says nothing favorable about
the candidate. The report must therefore store resolution status separately from
processing status; a single Boolean \texttt{done} is inadequate.

Materiality is another separate field. An issue is material when its unresolved
answer could change applicability, a required action, the decision, the timing,
the explanation required for review, or the validity of the evidence chain. If
impact is unknown, it is treated as material until assessed. Downgrading an issue
requires a reason and authority appropriate to the configured review policy.
A generator cannot label its own objection immaterial simply to unblock release.

A resolved issue can be reopened by new evidence, a source amendment, an invalid
witness or a discovered dependency. Reopening creates a new event and invalidates
results that depended on the old resolution. It does not erase the earlier
record. If the original run is already closed, the successor investigation links
to it and states which conclusion is being reconsidered.
% END SOURCE UNIT M06:07

% BEGIN SOURCE UNIT M06:08
\section{Choose a remedy that can answer the question}
\label{merge:M06:08}


Each discrepancy type has a permitted action set. A missing source calls for
retrieval or escalation to the source custodian. A parsing disagreement calls for
inspection of the original page, an alternative parser or manual transcription
with provenance. A scope disagreement calls for analysis of the actor and
activity language, incorporated definitions and the selected bank profile. A
formal mismatch calls for a replayable witness and a correction to the relevant
encoding or runtime component.

A useful action has a question, an expected evidence type, a bounded execution
method and a success condition. ``Ask another model'' is incomplete. ``Ask the
alternative reader to identify source-supported reasons why the exception might
attach to the whole offer, using only the retained paragraph and definition'' is
a bounded interpretive task. Its result is still a proposal, but the task makes
clear what new information it seeks.

The controller should prefer actions with a plausible path to resolution. If the
question is whether a fee was paid, querying transaction evidence may be more
useful than sampling another legal explanation. If the question concerns the
meaning of an incorporated term, more transaction data may be irrelevant. This
matching of remedy to issue is where structured uncertainty becomes operationally
valuable.

An action can also deliberately seek a falsifier. For a proposed rule, request a
feasible scenario in which it contradicts an express source condition. For a
classification, request a fact pattern that the definition treats counterintuitively
and ask which source justifies that boundary. For a claimed equivalence, ask the
solver for a difference rather than asking a model whether the programs look
similar. Falsification attempts must retain their domain and assumptions.

Repeated actions are allowed only when their rationale changes or a bounded retry
policy permits recovery from a technical failure. Retrying an unavailable server
can make sense; asking the identical question with no new evidence is unlikely
to resolve a legal ambiguity. The report should show which retries addressed
transport failure and which explored a substantive alternative. Both consume
resources, but they mean different things for the investigation.

A repair is validated against the original discrepancy and against the unaffected
requirements. Fixing the type-of-product omission must not remove the fee-discount
exception. A local witness becoming consistent is necessary evidence about that
repair, but insufficient evidence about the whole rule. The controller runs the
dependent validation set, records any new discrepancies, and keeps the old issue
open until its resolution criterion is actually satisfied.
% END SOURCE UNIT M06:08

% BEGIN SOURCE UNIT M06:09
\section{A resolution matrix prevents the wrong kind of closure}
\label{merge:M06:09}


The controller needs a typed relationship between a discrepancy and the evidence
that can close it. Without that relationship, every issue can eventually be
``resolved'' by another fluent explanation. The matrix below is a proposed
engineering contract for the first implementation. It states the minimum kind of
evidence, not a substitute for legal judgment about whether that evidence is
persuasive in the particular case.

\begin{longtable}{@{}P{27mm}P{56mm}P{65mm}@{}}
\caption{Discrepancy-specific actions and closure evidence.}\label{tab:resolution-matrix}\\
\toprule
Discrepancy & Appropriate next action & What can justify closure\\\midrule\endfirsthead
\toprule Discrepancy & Appropriate next action & What can justify closure\\\midrule\endhead
Missing source unit & Compare original rendering and independent inventory & Retained unit, corrected inventory and reviewed disposition; agreement among summaries is insufficient.\\
Unavailable dependency & Retrieve the named version or seek authorized clarification & Inspected applicable source, or an explicit adjudication of the remaining question; a guessed definition is insufficient.\\
Different source versions & Establish the assessment cut-off and amendment relationship & A documented version selection and reconsideration of affected candidates.\\
Actor or activity scope & Inspect addressee, operative language and profile facts & Reviewed applicability proposition with evidence for the actual bank activity.\\
Disputed classification & Identify definition and distinguishing factual evidence & An accepted definition and sufficient facts, or an attributable adjudication; a Boolean label alone is insufficient.\\
Exception attachment & Construct component and whole-arrangement examples & Source-grounded structural choice, with contrary reading disposition and regression cases.\\
Necessary versus sufficient condition & Compare implication directions on boundary scenarios & A reviewed direction of dependence and a rule that preserves other required conditions.\\
Timing or transition & Separate publication, commencement, valid time and knowledge time & Reviewed interval and transition rule with endpoint examples.\\
Factual conflict & Obtain correction or apply an approved evidence-priority rule & Attributable supersession or reviewed normalization; newest arrival alone is insufficient.\\
Formal mismatch & Produce and replay a feasible difference witness & Validated repair or domain-qualified equivalence, with source assumptions retained.\\
Unsupported operation & Inspect whether a supported encoding preserves meaning & Reviewed encoding plus conformance evidence, or an approved runtime extension; compiler rejection does not defeat the reading.\\
Java mismatch & Compare full semantic results and trace & Corrected exact package with repeated affected verification and new build identity.\\
Missing ensemble member & Retry under budget or use an authorized replacement & Required role output with documented independence limits; silence is not agreement.\\
Search truncation & Investigate retained frontier or narrow the claimed scope & Completed declared finite search or explicit reviewed restriction; a high best score is insufficient.\\
Broken approval lineage & Reconstruct exact source, report, bundle and build links & Fresh valid approvals for exact identities; a matching display title is insufficient.\\
\bottomrule
\end{longtable}

Some issues need more than one row. A cashback classification can depend on an
unavailable definition and an uncertain factual relationship to a fee. Retrieving
the definition resolves the dependency but may leave the factual question open.
The controller should represent these as linked issues so that progress is visible
without overstating closure. A single combined issue can be acceptable if its
resolution criterion explicitly requires both conditions.

The matrix also guides action rejection. If a proposed action cannot produce the
needed evidence, the scheduler should normally decline it and explain why. A new
model paraphrase is not an appropriate remedy for a corrupted source digest. A
solver equivalence query is not an appropriate remedy for a missing commencement
date. This does not prohibit exploratory work; it requires that exploratory work
be labelled as such and charged to the investigation budget rather than counted
as resolution.

Closure is evaluated against the issue's original question and current source
context. Suppose a model responds to a component-level question with a correct
whole-offer explanation. The response may be useful, but it has not answered the
original question. The controller creates or links the relevant candidate and
keeps the original issue open. This guards against a subtle form of apparent
progress in which the task changes until the model can answer it confidently.
% END SOURCE UNIT M06:09

% BEGIN SOURCE UNIT M06:10
\section{Resolution can reveal that the first question was wrong}
\label{merge:M06:10}


A structured controller must still permit substantive reframing. The initial
question may assume a distinction that the source does not make. For example,
the team may ask whether a cashback is a discount when the relevant provision
instead classifies the entire arrangement through another defined concept.
The correct response is to record the changed understanding, identify which
source supports it and replace the candidate with a new revision. The old question
is dispositioned as superseded by a reasoned reframing, not silently deleted.

Reframing does not reset budgets. The successor issue links to the original root
or consumes the same global account, and the report explains the relationship.
If the reframing expands the source perimeter beyond the authorized run, the
current investigation can end with a request for a successor run. This prevents a
useful ability to change direction from becoming an unbounded search loophole.

A concrete example concerns subject identity. A provider initially treats an
offer as one benefit. The independent inventory identifies a source passage that
requires separate treatment of components. The controller creates a component
model, revises factual requirements and reruns affected scenarios. It also records
that prior results used a different subject abstraction. Merely changing the
formula while leaving the whole-offer adapter untouched would not complete the
repair.

Another example concerns modality. A candidate initially turns guidance about
factors to consider into an unconditional prohibition. Investigation shows that
the source requires an assessment rather than a predetermined outcome. The repair
may need a workflow obligation, evidence of consideration and a review record,
not another Boolean threshold. If the current runtime cannot express the accepted
requirement, the result is a supported interpretation with an implementation gap.
That is more accurate than forcing the source into a familiar Boolean shape.

These cases explain why the issue register includes the affected decision and the
kind of evidence required. The controller can recognize that the target has
changed and invalidate dependent checks. A test suite for the previous target
cannot automatically certify the new one, even if many examples happen to retain
the same final status.
% END SOURCE UNIT M06:10

% BEGIN SOURCE UNIT M06:11
\section{The investigation graph should preserve shared causes}
\label{merge:M06:11}


The following diagram separates the objects that are often collapsed into a
single chain of model messages. Several candidates can depend on one definition,
and several discrepancies can affect one candidate. Actions investigate issues;
results create evidence or revisions. The final report reconciles the graph rather
than selecting the last message as the answer.

\begin{figure}[htbp]
\centering
\begin{tikzpicture}[
 box/.style={draw=teal,rounded corners,align=center,text width=43mm,minimum height=12mm,font=\small},
 arr/.style={-{Latex[length=2mm]},thick,draw=navy},node distance=12mm and 10mm]
\node[box] (source) {Retained source and\\reviewed inventory};
\node[box,right=of source] (definition) {Shared definition\\or missing premise};
\node[box,below=of source] (candidates) {Immutable competing\\candidate revisions};
\node[box,right=of candidates] (issues) {Discrepancies and\\issue-root budgets};
\node[box,below=of candidates] (evidence) {Validated evidence,\\witnesses and checks};
\node[box,right=of evidence] (actions) {Reserved bounded\\resolution actions};
\node[box,below=of evidence,xshift=26mm,text width=72mm] (report) {Complete report: resolved reasons,\\surviving alternatives and uncertainty};
\draw[arr] (source)--(candidates);
\draw[arr] (source)--(definition);
\draw[arr] (definition)--(issues);
\draw[arr] (definition)--(candidates);
\draw[arr] (candidates)--(issues);
\draw[arr] (issues)--(actions);
\draw[arr] (actions)--(evidence);
\draw[arr] (evidence)--(candidates);
\draw[arr] (evidence)--(report);
\draw[arr] (issues.east)--++(8mm,0)|-(report.east);
\end{tikzpicture}
\caption{Search creates a history, while dependencies explain what must be reconsidered.}
\label{fig:ensemble-dependencies}
\end{figure}

An important consequence is that one new source can resolve several issues
without spending a separate retrieval action for each. The action is charged once,
and each affected issue receives a validated link to the result. Their root
budgets do not acquire extra credits. Conversely, one failed source acquisition
can block several candidates, and the report should identify that shared cause
rather than presenting a long list of apparently unrelated uncertainties.

The graph also prevents misleading independence claims. If four candidate members
all rely on the same extracted definition and factual mapping, the report can
show that common dependency. Their reasoning paths may differ, but their evidence
base is shared. This is a more useful account of redundancy than simply displaying
four model logos beside an agreement percentage.
% END SOURCE UNIT M06:11

% BEGIN SOURCE UNIT M06:12
\section{Finite budgets are part of the semantics}
\label{merge:M06:12}


The user requires automatic processing up to a maximum number of times. To make
that requirement implementable, we distinguish an initial proposal phase from
resolution rounds. A resolution round consists of selecting currently open
issues, reserving authorized actions, collecting bounded results and reconciling
the resulting candidates. The round count increases when the controller commits
the round's action plan, even if every external action later times out.

The configuration supplies a maximum number of resolution rounds, a maximum
number of issued actions per issue root, a maximum number of issued actions for
the run, and an absolute deadline. It also supplies per-action timeouts, candidate
and graph-size limits, and any token or monetary limits. None is inferred from
the fact that the existing drafting stub happens to use three attempts. That
stub's retry behavior addresses a different contract.

A reference policy for deterministic fixtures uses three resolution rounds,
three issued actions per issue root, and twelve external actions in total,
including the initial proposal actions. It permits at most four initial external
proposal actions, leaving at most eight for later investigation. These are
engineering examples chosen to make boundary tests tractable. They are not
empirically justified production defaults or recommendations about how much
legal review is enough. Production configuration must be approved and evaluated
for the selected task and cost constraints.

The reference policy also sets a ten-minute run deadline and a thirty-second
external action timeout. Deterministic local checks have separate finite input
and execution bounds. A mandatory check is never silently skipped when its bound
is reached: it returns an incomplete result and blocks the corresponding claim.
Token and monetary caps are nullable only when the provider integration cannot
supply a reliable bound; such a configuration is labelled unsuitable for an
unattended paid run. A paid deployment must supply conservative reservations.

Let $B$ be the configured global action cap and $a$ the number of actions already
issued. Before dispatching any new external action, the controller must establish
\begin{equation}
 a+1\leq B,\qquad a_r+1\leq B_r,\qquad t<t_{\rm deadline},
 \label{eq:budget-guard}
\end{equation}
where $a_r$ and $B_r$ are the issued count and cap for the affected issue root.
For an initial proposal action, a separate initial-phase account replaces the
issue-root account. The guard also checks the round, token, cost and action-type
constraints. Passing one limit does not override another.

Reservation occurs atomically before dispatch. The issued count increases even if
the worker crashes before learning whether the provider accepted the request.
This conservative accounting prevents an ambiguous network result from becoming
free work. A reservation can be cancelled without spending only before dispatch
is durably committed and only through an atomic state transition that proves no
worker can still send the action.

An action's effective waiting limit is the smaller of its configured timeout
and the time remaining before the run deadline. The worker uses a monotonic clock
for elapsed waiting, while the durable UTC deadline supports restart and audit.
Clock adjustments cannot replenish an action budget. At the run deadline, pending
actions receive terminal or remote-result-unknown dispositions; the controller
does not extend the run simply because an action began a moment earlier.

Provider SDK retries must be disabled or included in the same accounting. A
single logical action that secretly makes five billable requests defeats the
meaning of the cap. When a provider offers idempotency keys, the adapter uses the
recorded action identifier. When it does not, the system cannot promise exactly
once execution across a lost response; it reports the uncertainty, counts the
attempt, and follows the configured retry policy.
% END SOURCE UNIT M06:12

% BEGIN SOURCE UNIT M06:13
\section{No progress is a substantive stopping reason}
\label{merge:M06:13}


A finite global budget guarantees that external actions cannot continue forever,
but it need not be spent in full. A no-progress condition avoids cycling through
identical proposals. The controller computes a progress signature from accepted
source additions, newly supported assumptions, new feasible witnesses, issue
resolutions and genuinely distinct candidate semantics. Cosmetic wording changes
do not count as progress.

The signature is a deterministic summary of validated changes, not a model's
claim that its answer improved. If a resolution round produces no such change,
the controller records that fact. The reference policy stops after two successive
rounds without progress, unless an already reserved action is still awaiting a
result within the deadline. This is an engineering stopping rule for the fixture
profile. A production policy may choose another threshold, but must disclose it
and evaluate the cost of premature stopping.

New unsupported candidates are not automatically progress. Otherwise a generator
could prevent termination by inventing fresh interpretations indefinitely.
Similarly, splitting one issue into several differently worded issues does not
reset the no-progress count or any budget. The controller retains those proposals
for review when useful, but their existence alone does not establish a reason to
continue automatic investigation.

A source becoming temporarily unavailable may justify a scheduled successor run,
but not an endless loop in the current run. The report records the failed
retrieval, the affected question and the conditions under which a new run would
be useful. A human or an authorized source-change service can initiate that run
with a new recorded budget. This preserves the distinction between bounded
automation and an operational process that may legitimately continue over days.

No-progress termination is not a judgment that the research direction has failed.
It says that the current actions, sources and budget have not resolved the
question. A later phase designed to obtain independent authority may still be
justified. The system must report whether it encountered an invalid source,
broken implementation, insufficient evidence, or simply a candidate that failed
its checks. Those outcomes call for different next steps.
% END SOURCE UNIT M06:13

% BEGIN SOURCE UNIT M06:14
\section{A deterministic first scheduler}
\label{merge:M06:14}


The first implementation needs a concrete action order before a learned search
policy is available. The reference scheduler uses a lexicographic key rather
than an opaque weighted score. It first handles integrity and source-coverage
problems, then missing dependencies and applicability, then definitions and
exception structure, then formal and Java mismatches, and finally report or
presentation discrepancies. This order is a proposed engineering policy: earlier
questions can invalidate the inputs to later checks. It is not a legal ranking
of every possible issue.

Within a class, unresolved material or unknown-materiality issues precede reviewed
immaterial issues. Earlier creation sequence precedes later sequence, and the
stable issue identifier breaks any remaining tie. Every selected action must be
permitted for the issue kind and have remaining root and global budget. If an
earlier issue is awaiting evidence that the available tools cannot obtain, it
remains visible while the scheduler can make progress on another issue. Priority
is not permission to spend the whole budget repeating an unavailable action.

For each issue, action selection follows a short decision procedure. Check whether
an already retained source or result answers the question; if so, validate and
link it without issuing a new external request. Otherwise, retrieve a named
missing authoritative dependency when a permitted retrieval route exists. If the
source is present but the candidate encoding differs, request a witness or a
source-grounded repair. If competing supported readings remain, select a
contrasting scenario or focused argument question. If no permitted action can
supply the required evidence, prepare the human review question and stop automatic
work on that issue.

Actions are not dispatched merely because the scheduler can name them. A
pre-dispatch validator checks that their inputs differ meaningfully from a prior
failed substantive action, unless the retry is an explicitly permitted transport
recovery. It also checks source identity, adapter support and the remaining
budget. Rejection is recorded as a scheduling disposition; it does not consume
an external-action slot unless a dispatch reservation has already been committed.
Local scheduling and validation remain bounded by graph and execution limits.

Candidate breadth uses a similarly explicit rule. At initial reconciliation,
reserve one retained slot for each supported interpretation family before
allocating additional slots to refinements. If the number of families exceeds the
candidate cap, retain the families in the same declared issue-priority and
creation order, store the excluded family identifiers in the frontier, and open
a search-incomplete issue. This policy favors reproducibility and visible
incompleteness. It does not claim that its ordering finds the best legal theory.

For the optional pair-disagreement heuristic, the first prototype uses equal
weights in equation~\eqref{eq:disagreement-question}. It enumerates the available
feasible contrasting scenarios within the supported comparison fragment, selects
the scenario separating the most retained pairs, and breaks ties by stable
scenario identifier. A cost-aware or learned weight can later replace this rule
under a versioned evaluation. Equal weights here are a baseline convenience,
not equal probabilities that every reading is correct.

The implementation report records the scheduler version and each selection key.
An engineer can therefore reconstruct why an issue was investigated before
another, and a reviewer can challenge the priority policy without guessing at a
model's hidden preference. Chapter~\ref{ch:evaluation} compares this deterministic
baseline with enhanced search at equal resources. Both must obey the same budget,
reporting and release constraints.
% END SOURCE UNIT M06:14

% BEGIN SOURCE UNIT M06:15
\section{A controller that can be implemented and checked}
\label{merge:M06:15}


The controller is a deterministic state machine around fallible external tools.
Its durable state contains the run configuration, source snapshot, required role
results, candidate revisions, issue roots, action records, budgets and report
status. External results enter through validation. They do not directly mutate
approvals or release state.

The following pseudocode specifies the main control flow. The functions named
\texttt{commit} and \texttt{reserve} are transactions; they must either write all
of their stated changes or none. The action worker uses the committed action
record and cannot create additional actions on its own.

\begin{lstlisting}[caption={Main controller, with explicit terminal uncertainty.}]
start_run(source_packet, profile, required_policy):
    validate_source_identity_and_policy()
    commit(run = INITIALIZING, counters = 0)
    reserve_and_issue_initial_roles_with_global_budget()
    collect_until_complete_or_action_deadline()
    validate_outputs_and_freeze_initial_proposals()
    reconcile_inventory_candidates_and_required_roles()

    while has_unresolved_issue():
        if integrity_failure():
            close_as(FAILED_INTEGRITY); break
        if cancelled():
            close_as(CANCELLED); break
        if deadline_or_global_or_round_limit_reached():
            mark_remaining_issues_with_stop_reason()
            close_as(BLOCKED_UNRESOLVED); break
        if no_progress_limit_reached():
            mark_remaining_issues(NO_PROGRESS)
            close_as(BLOCKED_UNRESOLVED); break

        plan = choose_permitted_discriminating_actions()
        if plan.is_empty():
            mark_remaining_issues(HUMAN_REQUIRED)
            close_as(BLOCKED_UNRESOLVED); break
        commit_next_round_and_reserve_bounded_actions(plan)
        dispatch_only_reserved_actions()
        collect_results_until_each_action_is_terminal()
        validate_results_and_create_candidate_revisions()
        rerun_affected_checks_and_reconcile_all_open_issues()
        record_progress_signature()

    if not already_closed():
        apply_integrity_cancellation_and_deadline_guards()
    if not already_closed():
        check_required_roles_coverage_and_all_mandatory_checks()
        if any_requirement_incomplete():
            close_as(BLOCKED_UNRESOLVED)
        else:
            close_as(READY_FOR_REVIEW)
    write_complete_report_for_every_terminal_status()
    return report_reference
\end{lstlisting}

Every detected discrepancy enters reconciliation, including differences later
judged immaterial. A deterministic presentation-equivalence check can close a
purely cosmetic issue immediately with its reason; unassessed differences remain
unresolved and conservatively material. The loop therefore processes unresolved
issues rather than only issues already assigned a high severity. Integrity,
cancellation and deadline guards are checked again before final readiness, even
when no issue remains, so an empty loop cannot bypass a terminal condition.

The final coverage check is essential. If no comparator detects a disagreement,
the loop might never run. The controller must still require complete source
coverage, mandatory members and verification results. ``No open issue'' and
``all required evidence present'' are separate conditions. A failed member must
also create an issue, but the final check protects against a defect in that issue
creation path.

Actions that finish after the run deadline are retained as late observations.
They do not reopen a closed run or retroactively change its result. A successor
run may use them after validation. This rule makes reports stable and prevents a
late worker from changing the package that a reviewer has already inspected.
Late evidence can still trigger a documented follow-up when it changes risk.

The controller must produce a report even when it fails. If the primary report
writer is unavailable, the durable run and action records support recovery of a
minimal failure report. An integrity failure prevents a positive result, but it
must not erase the explanation of what was attempted. Report generation itself
has bounded resource use and a recovery path; an exception thrown while rendering
a large argument graph cannot make the investigation disappear.
% END SOURCE UNIT M06:15

% BEGIN SOURCE UNIT M06:16
\section{Termination and the boundary of the proof}
\label{merge:M06:16}


We can prove a limited but useful property of the proposed controller. Assume a
finite input packet, a finite configured action cap $B$, atomic issue accounting,
finite local reconciliation with explicit graph limits, and an enforced timeout
for every dispatched action. Assume also that the storage and supervisor either
respond within their operational bounds or move the run to a recoverable failure
state. These assumptions concern implementation and infrastructure; the proof
does not establish them merely by listing them.

Define the remaining action budget as $V=B-a$, where $a$ is the durable number of
issued actions. Initially $V$ is a nonnegative integer. Each new dispatch requires
an atomic reservation and reduces $V$ by one. No transition increases $V$ in the
same run. Therefore there can be at most $B$ dispatches. Between dispatches,
reconciliation is finite by its input and execution bounds. Waiting is finite by
the action timeout and absolute deadline. When no permitted action remains, the
controller enters a terminal state. Under these assumptions, automatic processing
terminates.

This argument remains valid when an action discovers additional issues. New
issues do not replenish the global budget. It also remains valid when a worker
restarts, provided the durable counter is authoritative and a lease does not
authorize a second unrecorded dispatch. The per-issue cap and round cap are useful
additional constraints, but the global cap is what prevents an unbounded stream
of newly named issues from evading termination.

The proof is about bounded processing, not resolution. A controller can terminate
with every difficult issue unresolved. It is also not a proof of wall-clock
availability under arbitrary storage failure. If the process is indefinitely
suspended or the database never responds, no local algorithm can guarantee a
report at a fixed real-world time without an external supervisor. The operational
design must therefore monitor leases and deadlines independently and recover a
failure report when infrastructure returns.

A second invariant concerns release. Let $M(r)$ mean that run $r$ has an unresolved
material issue, $C(r)$ mean that its mandatory coverage and checks are complete,
and $A(r,b)$ mean that the required authorized reviewers approved exact bundle
$b$ on the evidence from $r$. The release service requires
\begin{equation}
 \operatorname{Release}(r,b)\Longrightarrow
 \neg M(r)\land C(r)\land A(r,b)\land\operatorname{Integrity}(r,b).
 \label{eq:ensemble-release}
\end{equation}
This is a design invariant to be enforced by the release transaction. The
controller's \texttt{READY\_FOR\_REVIEW} state supplies neither $A$ nor release
itself. An implementation test must attempt to bypass each conjunct, including
through a stale approval or a source update, and verify rejection.

The invariant also exposes the limit of the design. If the inventory misses a
material source and no check detects it, $C(r)$ could be incorrectly recorded.
Formal verification of the release transaction would still prove only that it
uses the recorded evidence correctly. Independent source review and empirical
omission tests address the reliability of that evidence. The book's separation
of proof claims is therefore retained inside the ensemble rather than abandoned
at its most important decision.
% END SOURCE UNIT M06:16

% BEGIN SOURCE UNIT M06:17
\section{Durable actions, leases and ambiguous execution}
\label{merge:M06:17}


An action moves through \texttt{RESERVED}, \texttt{DISPATCHED}, and one of
\texttt{SUCCEEDED}, \texttt{FAILED}, \texttt{TIMED\_OUT}, or
\texttt{RESULT\_UNKNOWN}. The last state is necessary when the controller cannot
determine whether a remote service performed the action. A lost connection after
submission is different from a request rejected before dispatch. Both can lead
to missing results, but only the former may already have consumed remote work.

The reservation transaction records the action identifier, input hashes, issue
root, round, budget charge, adapter identity, timeout and idempotency key. An
outbox entry is written in the same transaction. A worker claims that entry with
a lease and a fencing token. The token increases when ownership changes, so an
old worker cannot commit a result after a newer owner has taken responsibility.
A valid late result can be retained as an observation without changing the closed
action's authoritative state.

The dispatch boundary requires care. If the outbox records only ``not yet sent''
and a worker crashes after sending but before updating the database, another
worker may send again. A provider idempotency key can make that retry refer to the
same operation, but not every provider supports this guarantee. The adapter must
state its capability. Without remote idempotency, a recovery policy can choose to
leave the result unknown rather than risk another paid request, or authorize a
new counted action. It cannot truthfully promise exactly once remote execution.

The budget system reserves an upper bound on tokens or cost when those resources
are capped. Settlement records actual usage when known. A refund of unused
monetary reservation may be legitimate after a confirmed terminal response, but
it does not decrement the issued-action count. Otherwise a sequence of failed
requests could evade the maximum number of attempts. If actual usage exceeds a
provider's documented bound, the run records an accounting discrepancy and stops
new spending until the integration is investigated.

A compare-and-swap version protects issue updates. The worker returns evidence
against the issue revision it received. If the issue has changed, the controller
revalidates whether the evidence is still relevant before applying a resolution.
This prevents two simultaneous actions from overwriting each other's objections
or resolving an issue using a candidate that is no longer current.

Cancellation stops new reservations and attempts to cancel outstanding work.
Remote cancellation may be best effort. The report records outstanding or late
actions and keeps their budget charges. Cancellation cannot convert an incomplete
run into a clean result, and it cannot delete the source of a later adverse
finding. These details matter because regulatory evidence must survive ordinary
operational events, including a person closing a browser or restarting a worker.
% END SOURCE UNIT M06:17

% BEGIN SOURCE UNIT M06:18
\section{What counts as a justified repair}
\label{merge:M06:18}


A repair has two responsibilities: remove the identified defect and preserve the
requirements that were already correct. For a missing product-type condition,
the revised body must include the type route, and the previously verified gift,
exception and scope behavior must remain intact. The repair record identifies the
changed expression and the affected checks. It cannot simply report that a new
model answer received a higher score.

The original discrepancy remains immutable. The repair links to it and supplies
resolution evidence: the source clause, changed candidate, witness before and
after, regression checks and any reviewer decision. A failed repair becomes
another candidate revision with a failed status. This is useful evidence about
the investigation and should not be removed from the dossier merely because a
later revision passes.

A repair can be rejected at several layers. The source may not support its new
condition. The controlled-language rendering may no longer match the formal
rule. The formal rule may use an unsupported operation. The Java package may
miscompute the accepted expression. These are different failures. A compiler
rejection means the current encoding cannot be delivered through that compiler;
it does not necessarily refute the legal reading. The next action might be a
supported encoding or a documented implementation extension rather than a new
interpretation.

Reference cases require special protection. When a candidate disagrees with a
frozen expected outcome, the system records the disagreement and investigates
both sides. It does not assume the reference is infallible, particularly in our
same-author 22-case exercise. But changing the reference requires an independent
resolution record with evidence and reviewer authority. A model cannot rewrite
the expected result and then claim that its unchanged candidate now passes.

The same rule applies to formal domains. If a witness exposes a defect, narrowing
the domain to exclude it is a substantive change. The domain change needs a
source or operational justification, not merely a desire to eliminate the
counterexample. Every equivalence claim includes the domain hash and assumptions,
so an excluded scenario remains visible to reviewers.

Finally, repair must preserve dissent. If one supported reading is selected after
review, the rejected reading, its strongest argument and the reason for rejection
remain attached. Future evidence may make that earlier objection relevant. A
system that stores only the winning candidate forces the next team to rediscover
the same ambiguity and cannot explain why the original team chose as it did.
% END SOURCE UNIT M06:18

% BEGIN SOURCE UNIT M06:19
\section{Scores that assist investigation without pretending to be probabilities}
\label{merge:M06:19}


A single confidence number combines incompatible questions. A source can be
authoritative while its application is ambiguous. A candidate can be internally
consistent while missing a clause. A Java artifact can be verified against an
unapproved rule. The proposed output therefore uses a score vector and explicit
statuses rather than a scalar legal-confidence percentage.

The vector records source coverage, dependency completeness, supported and
unsupported assumptions, unresolved material objections, scenario agreement,
formal-check status, implementation-preservation status and search coverage.
Some components are counts or categorical states. Others may be empirical
scores, such as retrieval relevance. Their definitions are stored with the run.
A score is not allowed to compensate for a failed mandatory condition.

For example, a candidate might cover eighteen of eighteen inventoried paragraphs
with dispositions, yet have two unresolved incorporated sources and one material
classification question. Displaying only ``100 percent coverage'' would mislead.
The report must state that source-unit disposition is complete while dependency
and interpretation resolution are incomplete. Conversely, a lower model ranking
need not disqualify a candidate whose source argument remains undefeated.

Search priority can be a weighted function of expected information gain, action
cost and issue materiality. The weights are scheduling choices. They may be tuned
under a declared evaluation protocol, but they are not probabilities of legal
correctness. Every displayed numerical score includes its purpose and calibration
status. A field named \texttt{probability\_correct} remains absent or null until an
appropriate target-specific calibration study supports a precisely defined event.

A useful shortlist can be a Pareto set. Candidate $h$ dominates candidate $j$ only
when it is no worse on every declared comparable dimension and better on at least
one. Even then, dominance is an aid to presentation, not legal defeat. A candidate
with a unique source argument should not disappear merely because another has
more completed engineering checks. Unsupported candidates can be rejected for
specific reasons; incomparable supported candidates remain visible.

Set-valued reporting is especially valuable when authority is genuinely
incomplete. The dossier may present two conditional proposals: under definition
$D_1$, use rule $R_1$; under definition $D_2$, use rule $R_2$. It must identify who
can resolve the definition and what operations remain blocked. This is more
informative than selecting one proposal with a decorative confidence interval
whose coverage has never been established.
% END SOURCE UNIT M06:19

% BEGIN SOURCE UNIT M06:20
\section{A complete dry run: repairing an explicit omission}
\label{merge:M06:20}


The following trace is a designed example for the future controller. It has not
been executed by a live model and does not replace the retained second-circular
results. The source packet contains the selected paragraph, its surrounding
context and the unresolved dependency records. The activity profile and factual
schema are fixed before proposals begin.

Four initial actions are reserved. The normative reader proposes
$G\land P\land\neg D$, omitting product-type linkage. The controlled-language
reader proposes $G\land(P\lor T)\land\neg D$. The alternative reader agrees that
both product and type links require examination and identifies a separate
classification ambiguity. The inventory member confirms the paragraph and
footnote enumeration but reports an unavailable incorporated source. The run has
issued four actions; eight remain under the reference global cap.

Reconciliation creates issue $I_1$ for the missing type route, $I_2$ for the
classification ambiguity and $I_3$ for the unavailable dependency. The comparator
does not average the two formulas. It constructs the known-fact scenario
$G=\True$, $P=\False$, $T=\True$, $D=\False$, with scope established. The first
candidate returns false and the second true. The source passage explicitly
includes the type route, so $I_1$ has a concrete textual and behavioral basis.

In resolution round one, the controller reserves a source-grounded repair action
for $I_1$ and a dependency retrieval action for $I_3$. The issued count becomes
six before either is dispatched. The repair produces a child candidate that adds
$T$ to the disjunction. Its controlled-language rendering, type checks and witness
replay pass. The regression set checks that fee discounts remain excluded and
scope uncertainty retains its established behavior. The issue is resolved by
evidence with the original discrepancy and both candidate versions retained.

The dependency retrieval times out. Its action is terminal and counted; $I_3$
remains unresolved. The run does not claim complete circular interpretation just
because $I_1$ was repaired. If the unavailable source could affect the selected
classification, both $I_2$ and the relevant candidate remain blocked. The report
can still state precisely that the explicit type-of-product omission was found
and corrected.

This trace illustrates two kinds of progress. The repaired omission supplies
positive evidence about the ensemble's ability to catch that seeded defect in
this designed case. The timeout supplies operational evidence about the
controller's accounting. Neither is a measured error rate on real circulars.
The prototype acceptance test will reproduce these transitions with scripted
members before any live-model comparison is attempted.
% END SOURCE UNIT M06:20

% BEGIN SOURCE UNIT M06:21
\section{A complete dry run: an ambiguity that survives the budget}
\label{merge:M06:21}


Continue with the hypothetical cashback question. One candidate treats a later
payment as a fee discount when it reimburses an actually paid charge and is capped
at that charge. Another requires a reduction in the contractual fee itself.
Their rules agree on a direct waiver and disagree on a later reimbursement.
Neither definition is accepted merely because it produces a convenient outcome
for the bank.

The scenario designer creates three synthetic offers. Offer A waives a charge of
HK\$100. Offer B charges HK\$100 and later pays HK\$100 back under a linked
promotion. Offer C charges HK\$100 and pays HK\$150. The factual amounts are known,
but the legal classification of the later payment remains disputed. The dossier
shows the candidates' conclusions and the exact definition each requires. It
also separates any excess payment into a proposed component analysis rather than
silently treating the entire amount as a discount.

In round one, a retrieval action seeks the incorporated definition. If the source
is inaccessible, the result is a failed retrieval, not support for either
reading. In round two, a bounded argument action asks each candidate to identify
the strongest source-based reason against its own definition. Suppose both
return plausible but unresolved explanations and no new authority. The controller
records the arguments, but a longer explanation alone does not constitute
resolution.

In round three, a question-selection action prepares a focused review question
with the contrasting offers and the missing premise. The action can successfully
produce the question while the issue remains unresolved. After the third round,
no further automatic resolution round is authorized. The candidate set is
preserved, the issue receives a budget stop reason, and the run closes as
\texttt{BLOCKED\_UNRESOLVED}.

The report states which offers are affected, what the candidates agree on, where
they differ, which sources were inspected, which source remains unavailable, and
what an authorized reviewer must decide. It contains no assertion that three
rounds make the leading reading probably correct. The unresolved classification
is not replaced by a false Boolean, and the Java release route remains blocked
for the affected control.

A later authorized reviewer may establish the definition from an applicable
source or decide that a conservative bank policy should govern the uncertain
category. Those are distinct decisions. A bank policy that declines all cashback
offers pending clarification can be operationally useful, but it must be labelled
as the bank's precautionary choice. It cannot be reported as a proven SFC
prohibition unless the source supports that proposition.
% END SOURCE UNIT M06:21

% BEGIN SOURCE UNIT M06:22
\section{The uncertainty report is a first-class result}
\label{merge:M06:22}


An uncertainty report must let a reviewer reconstruct the decision without
reading raw model transcripts. It starts with the source and activity perimeter,
the proposed control, the run's terminal status and the operational consequence.
A blocked report should say what is blocked. ``Low confidence'' is not enough:
the reader needs to know whether the uncertainty concerns an incentive category,
an effective date, a missing definition or the integrity of the entire source.

The report then presents the surviving interpretations in comparable form. For
each, give the controlled-language statement, formal rule identifier, assumptions,
supporting sources, strongest objections and distinguishing scenarios. Use the
same factual scenarios across candidates so that the difference is visible.
Candidate ordering may reflect a declared presentation policy, but must not imply
that the first row has been approved.

Every detected discrepancy appears in an issue register with its processing and
resolution states, materiality, affected decisions, attempted actions and final
reason. Resolved issues remain present because their history explains why the
final candidate differs from the initial one. Unresolved issues identify the
missing evidence or human decision. A report that shows only unresolved issues
would conceal the repairs; a report that shows only repaired issues would conceal
the remaining risk.

Coverage is reported against explicit denominators. Count source units inventoried
and dispositioned, dependencies required and acquired, interpretation families
identified and explored, mandatory roles completed, candidate comparisons
attempted, and verification obligations passed or incomplete. The denominators
must be retained even when the run exceeds a display limit. Large inventories
can be paginated or stored as linked machine-readable records; they cannot be
silently truncated to make the report look complete.

The execution account records configured limits, issued actions, rounds used,
timeouts, unknown remote results, no-progress decisions, elapsed time, token and
cost reservations, actual usage when available, and unspent budget. A run that
stops early because it lacks authority should not look like a run that exhausts
its resources. Both may be blocked, but their next useful actions differ.

A concise machine-readable summary might look as follows. Identifiers here are
illustrative; the complete contract in Chapter~\ref{ch:implementation} adds
validated references and integrity fields.

\begin{lstlisting}[caption={A blocked result preserves alternatives and the next decision.}]
{
  "schema_version": "interpretation.v1",
  "run_status": "BLOCKED_UNRESOLVED",
  "source_ref": "23EC46-retained-version",
  "selected_slice": "paragraph-10",
  "candidate_ids": ["cashback-reimbursement", "cashback-fee-reduction"],
  "material_unresolved_issue_ids": ["classification-cashback"],
  "release_eligible": false,
  "processing_stop": "ROUND_LIMIT",
  "rounds_issued": 3,
  "probability_of_legal_correctness": null,
  "next_decision": "Establish the applicable fee-discount definition",
  "next_decision_owner_role": "LEGAL_INTERPRETATION_REVIEWER"
}
\end{lstlisting}

The report should also state the limits of its evidence. In the second-circular
exercise, the expected cases and candidate were prepared by the same author.
That fact belongs beside the verification summary. A reviewer should not have to
find a footnote in a separate research note to discover that the apparent oracle
was not independent. Similarly, a solver timeout must appear as incomplete
verification, not disappear from an average pass rate.

Reports are immutable versions. A corrected report points to its predecessor and
explains the correction. Approval binds the exact report, candidate, source
packet and Java bundle. The report is thus both a readable explanation and an
index into reproducible evidence. Its human readability matters because a
technically complete JSON object that no reviewer can understand does not solve
the original manual-interpretation problem.
% END SOURCE UNIT M06:22

% BEGIN SOURCE UNIT M06:23
\section{Agreement can still conceal a missed obligation}
\label{merge:M06:23}


The hardest example is not a disagreement. Suppose every generator ignores a
footnote that changes applicability. All formulas, scenarios and Java results
agree because they share the same omission. Repeated self-consistency sampling
would reinforce that agreement. Roundtrip translation could reproduce the same
shortened meaning. Differential execution would confirm that the wrong
specification is implemented consistently.

The independent inventory creates another opportunity to catch the omission. It
requires a disposition for the footnote. A normative reader responsible for that
unit may identify the qualification even if the candidate generators did not.
A source-to-rule coverage check then opens an issue because the footnote has no
corresponding assumption or rule. The repair propagates through candidate
comparison and implementation checks.

Now suppose the inventory also misses the footnote because both extraction paths
use the same defective text layer. Rendered-document inspection or a distinct
acquisition route may detect it. If those are absent or also fail, the system can
still miss the obligation. This remaining possibility is why the product's
claims must be measured against independent source review and adversarial
omission tests. Redundancy reduces opportunities for undetected error only to
the extent that the redundant activities can fail differently.

A second common-mode failure concerns vocabulary. All members may agree that a
payment is a discount because the prompt labels it ``fee rebate.'' The remedy is
to distinguish observed facts from proposed classifications. Present the payment,
charge, contractual relationship and timing without inserting the disputed label
as an established input. Then ask what evidence supports the classification.
Blindness to the preferred answer is useful only if the factual packet itself
has not already embedded that answer.

A third failure concerns evaluation. If the reference reviewer sees the model's
answer before forming an expected outcome, the reference may inherit the model's
mistake. The ensemble may then appear accurate against a contaminated oracle.
Independent initial judgments and recorded adjudication are necessary to measure
whether the product actually discovers errors. They are also part of a credible
operational review process, especially for novel or high-impact controls.
% END SOURCE UNIT M06:23

% BEGIN SOURCE UNIT M06:24
\section{Interaction between uncertainty and operational action}
\label{merge:M06:24}


The interpretation service should not decide every operational response. Its job
is to state what is known, what is disputed and which controls are justified by
the approved specification. The bank's host system then applies an explicit
operational policy to outcomes such as unknown, conflict or blocked release.
Those policies may include stopping an offer, routing a case to review or using
an already approved fallback control.

This separation avoids a dangerous ambiguity in ``fail closed.'' Closing an
uncertain marketing offer may be appropriate in a particular process, but
stopping a payment, withholding a disclosure or disabling an existing control
could create other obligations or harms. The responsible business and control
owners must approve the action for the relevant workflow. The interpretation
engine supplies a reasoned status; it does not invent a universal response for
all banking operations.

For the selected gift rule, a true result indicates that the specified
prohibition condition is satisfied. A false result says only that this condition
is not established under the accepted facts and rule. It does not authorize the
offer under every other requirement. The host must combine this result with other
applicable controls, such as disclosure or suitability obligations, according to
an independently reviewed orchestration policy.

An unresolved interpretation can sometimes be isolated. If the disputed cashback
definition affects only one incentive category and the system can reliably
identify that category, a reviewer may approve a limited control for unaffected
cases while routing the disputed category to manual review. The scope restriction
must be explicit, tested and included in the approved bundle. It is not enough to
state informally that the system will ``avoid edge cases.''

If impact cannot be isolated, the block applies more broadly. The system should
prefer an honest incomplete release proposal over an unsupported claim of
separability. The uncertainty report gives the decision-makers the information
needed to choose between narrowing scope, obtaining more evidence, retaining an
existing control or postponing deployment.
% END SOURCE UNIT M06:24

% BEGIN SOURCE UNIT M06:25
\section{Acceptance conditions for the central method}
\label{merge:M06:25}


The first implementation should be accepted against deterministic scenarios
before live model performance is assessed. A scripted member can deliberately
omit the type route, supply a contradictory source version, fail to respond, or
return a malformed candidate. These fixtures isolate controller behavior from
model variability. Passing them establishes specified engineering behavior, not
legal translation accuracy.

The controller must demonstrate that a unanimous but uncovered source unit opens
an issue; that a supported minority candidate remains in the report; that a
formal mismatch produces a replayable witness; that an unavailable dependency
cannot be marked resolved by repeated explanation; and that a successful repair
reruns affected checks. Each condition targets a concrete way an apparently
successful ensemble could hide an error.

Budget tests must reach exact boundaries. With a cap of twelve, the thirteenth
external dispatch must be rejected even if it is proposed by a new child issue.
A crash after dispatch must not reset the count. A timeout must consume its
issued action. A new worker with an old lease must fail to overwrite the current
result. A provider retry must either reuse a verified idempotency key or consume
a separately authorized attempt.

Report tests must verify completeness for successful, blocked, cancelled and
failed runs. Every issue identifier must have a disposition; every unresolved
material issue must appear in the blocking summary; every candidate revision
must retain its parent and source references; and every displayed score must
identify its purpose. Empty extension sets, solver unknowns and missing member
results must remain explicit rather than being converted to success by default
logic.

The release test is adversarial. Attempt release from
\texttt{READY\_FOR\_REVIEW} without approvals, with an approval for a different
bundle, after a source change, with an unresolved material issue, and after an
integrity failure. All attempts must fail. A positive release fixture requires
complete checks and authenticated approvals for the exact package. The local
synthetic identity mechanism is sufficient to test the state transitions, but
enterprise authentication and role enforcement remain a deployment requirement.

These conditions give an implementation agent a concrete target. They also
preserve the intellectual purpose of the ensemble: investigate alternatives,
expose missing evidence, and retain uncertainty when the evidence does not settle
the question. The resulting product is useful because it makes responsibility
and reasoning visible, not because a collection of fallible components can be
renamed a proof of English meaning.
% END SOURCE UNIT M06:25

